%% RAMSES Data Structure Description for LaRT v2.00
%% Author: Kwang-Il Seon
\documentclass[english,a4paper]{kiseon_memo}
\synctex=1
\usepackage{listings}
\usepackage{xcolor}
\usepackage{booktabs}
\usepackage{multirow}

\lstdefinestyle{fortranstyle}{
  basicstyle=\ttfamily\small,
  keywordstyle=\color{blue}\bfseries,
  commentstyle=\color{gray}\itshape,
  stringstyle=\color{olive},
  frame=single,
  breaklines=true,
  numbers=left,
  numberstyle=\tiny\color{gray},
  language=[95]Fortran,
  morekeywords={read,write,open,close,allocate,deallocate}
}
\lstdefinestyle{shellstyle}{
  basicstyle=\ttfamily\small,
  frame=single,
  breaklines=true,
  language={}
}
\lstset{style=fortranstyle}

\begin{document}

\title{RAMSES Data Structure and the LaRT Interface}
\author{Kwang-Il Seon}
\date{2026-04-18}
\maketitle

\tableofcontents
\bigskip

%=============================================================================
\section{Overview of RAMSES}
%=============================================================================

RAMSES \cite{teyssier2002} is an adaptive mesh refinement (AMR) code for
astrophysical hydrodynamics and N-body dynamics.  It decomposes the
computational domain into an octree hierarchy of grids, refining regions of
interest to arbitrarily high resolution.  The MPI decomposition uses a
space-filling curve (Peano--Hilbert ordering) to distribute octs among CPU
ranks.

\texttt{LaRT~v2.00} can read RAMSES binary output directly via
\texttt{amr\_type = 'ramses'}.  This document describes the RAMSES binary
format in detail and explains how \texttt{read\_ramses\_amr.f90} translates
it into the flat leaf-cell arrays required by the LaRT octree engine.

%=============================================================================
\section{Output Directory Layout}
%=============================================================================

A RAMSES snapshot is stored in a directory named
\texttt{output\_NNNNN/} where \texttt{NNNNN} is the five-digit output number.
For a run with $N_\mathrm{cpu}$ MPI tasks the snapshot contains the following
per-CPU files:

\begin{lstlisting}[style=shellstyle]
output_00010/
    info_00010.txt            ! global metadata: units, cosmology, etc.
    amr_00010.out00001        ! AMR tree for CPU 1
    amr_00010.out00002        !   ...
    ...
    amr_00010.outNNNNN        ! AMR tree for CPU N
    hydro_00010.out00001      ! hydrodynamic variables for CPU 1
    hydro_00010.out00002
    ...
    hydro_00010.outNNNNN
    gravity_00010.out00001    ! gravitational potential (optional)
    part_00010.out00001       ! particle data (optional)
\end{lstlisting}

All \texttt{.out*} files are Fortran \emph{unformatted sequential} binary
files (each record bracketed by 4-byte integer record markers on little-endian
systems).

%=============================================================================
\section{The Info File (\texttt{info\_NNNNN.txt})}
%=============================================================================

\texttt{info\_NNNNN.txt} is a plain-text file containing global simulation
parameters.  The fields relevant to LaRT are listed below.

\begin{center}
\begin{tabular}{lll}
\toprule
Key & Type & Meaning \\
\midrule
\texttt{ncpu}    & integer & Number of MPI ranks used for this snapshot \\
\texttt{ndim}    & integer & Number of spatial dimensions (always 3 for 3-D runs)\\
\texttt{nlevelmax} & integer & Maximum AMR refinement level used \\
\texttt{boxlen}  & real & Box length in \emph{code units} (see below)\\
\texttt{unit\_l} & real & Length unit: 1 code length $=$ \texttt{unit\_l} cm \\
\texttt{unit\_d} & real & Density unit: 1 code density $=$ \texttt{unit\_d} g~cm$^{-3}$ \\
\texttt{unit\_t} & real & Time unit: 1 code time $=$ \texttt{unit\_t} s \\
\texttt{gamma}   & real & Adiabatic index $\gamma$ of the EOS \\
\bottomrule
\end{tabular}
\end{center}

The velocity unit is derived as $\texttt{unit\_v} = \texttt{unit\_l}/\texttt{unit\_t}$
(cm~s$^{-1}$).

\paragraph{Cosmological runs.}
For cosmological simulations, \texttt{boxlen} is in Mpc$/h$ and
\texttt{unit\_l} incorporates the scale factor $a$ so that
$L_\mathrm{phys} = \texttt{boxlen} \times \texttt{unit\_l}$~cm is the
comoving box size.  LaRT reads the physical box size as
$L_\mathrm{cm} = \texttt{boxlen} \times \texttt{unit\_l}$.

\paragraph{Example info file (non-cosmological).}
\begin{lstlisting}[style=shellstyle]
 ncpu        =           64
 ndim        =            3
 nlevelmax   =           12
 ngridmax    =       500000
 boxlen      =   1.0000000000000000E+02
 unit_l      =   3.0856775814913670E+21   ! 1 kpc in cm
 unit_d      =   6.7707722919716770E-25
 unit_t      =   3.0856775814913670E+14
 gamma       =   1.6666666666666667E+00
\end{lstlisting}

In this example the box is 100~kpc on a side.

%=============================================================================
\section{The AMR File (\texttt{amr\_NNNNN.outCCCCC})}
%=============================================================================

Each CPU file stores the AMR tree data for that CPU's domain.  The file
consists of a \emph{header section} followed by a sequence of
\emph{level blocks}, one per AMR level.

\subsection{Header Records}

The AMR header records are read in the following order:

\begin{center}
\begin{tabular}{rll}
\toprule
Record & Variable(s) & Description \\
\midrule
1  & \texttt{ncpu}          & Total number of MPI ranks \\
2  & \texttt{ndim}          & Number of dimensions \\
3  & \texttt{nx, ny, nz}    & Coarse grid size (typically $1\times1\times1$ or $2\times2\times2$) \\
4  & \texttt{nlevelmax}     & Maximum AMR level \\
5  & \texttt{ngridmax}      & Maximum number of grids (octs) per CPU \\
6  & \texttt{nboundary}     & Number of boundary domains \\
7  & \texttt{ngrid\_current}& Current number of active grids \\
8  & \texttt{boxlen}        & Box length (code units) \\
9--21 & (various)           & Cosmological parameters, time, etc.\ (skipped by LaRT)\\
22 & \texttt{ngridlevel(:,:)} & Grid count array: \texttt{(ncpu, nlevelmax)} \\
23 & (level sums)           & Skipped \\
24--27 & (boundary grids)   & Only present if \texttt{nboundary > 0} (skipped) \\
28--32 & (ordering arrays)  & Peano--Hilbert keys, son maps, etc.\ (skipped) \\
\bottomrule
\end{tabular}
\end{center}

\paragraph{The ngridlevel array.}
\texttt{ngridlevel(j, l)} gives the number of octs belonging to CPU $j$
at AMR level $l$.  LaRT constructs a combined table
\texttt{ngridfile(j, l)} that includes both domain and boundary grids.

\subsection{Level Blocks}

After the header, the file contains one block per AMR level
$l = 1, \ldots, \texttt{nlevelmax}$.  Within each level, data are stored
\emph{domain by domain}: first all domains $j = 1, \ldots, \texttt{ncpu}$,
then all boundary domains $j = \texttt{ncpu}+1, \ldots,
\texttt{ncpu}+\texttt{nboundary}$.

For each domain $j$ at level $l$ (if \texttt{ngridfile}$(j,l) > 0$):

\begin{center}
\begin{tabular}{rll}
\toprule
Record(s) & Variable & Description \\
\midrule
1     & \texttt{ind\_grid(:)}  & Global grid indices \\
2     & \texttt{next(:)}       & Linked-list next pointer \\
3     & \texttt{prev(:)}       & Linked-list previous pointer \\
4--6  & \texttt{xg(:, 1:3)}   & Oct centre coordinates in code units \\
7     & \texttt{father(:)}     & Parent grid index \\
8--13 & \texttt{nbor(:,1:6)}  & Neighbour grid indices (6 faces) \\
14--21 & \texttt{son(:,1:8)}  & Child grid indices; $=0$ if a child is a \textbf{leaf} \\
22--29 & \texttt{cpu\_map(:,1:8)} & CPU ownership per child \\
30--37 & \texttt{ref\_map(:,1:8)} & Refinement flags per child \\
\bottomrule
\end{tabular}
\end{center}

Record counts assume \texttt{ndim}$=3$, \texttt{twotondim}$=8$.

\subsection{Oct Structure and Leaf Identification}

RAMSES groups cells into \emph{octs}: cubes of $2^3 = 8$ cells sharing a
common parent grid cell.  Each oct is identified by its \emph{grid index}.
Within an oct, the eight child cells are indexed $\texttt{ind} = 1, \ldots, 8$
with the following ordering:

\begin{equation}
  i_x = (\texttt{ind}-1) \bmod 2, \quad
  i_y = \lfloor(\texttt{ind}-1)/2\rfloor \bmod 2, \quad
  i_z = \lfloor(\texttt{ind}-1)/4\rfloor,
\end{equation}
so that \texttt{ind}=1 is the ($-x,-y,-z$) corner and \texttt{ind}=8 is the
$(+x,+y,+z)$ corner.

A child cell at index \texttt{ind} in grid $g$ at level $l$ is a
\textbf{leaf} if and only if
\begin{equation}
  \texttt{son}(g, \texttt{ind}) = 0.
\end{equation}
Otherwise the value of \texttt{son} is the grid index of the finer oct that
occupies that child cell.

\subsection{Cell Position Calculation}

The centre of the coarse grid of the root domain is the origin.
For a non-cosmological run with \texttt{nx=ny=nz=1}:
\begin{equation}
  \mathbf{x}_\mathrm{bound} = \left(\frac{n_x}{2},\, \frac{n_y}{2},\, \frac{n_z}{2}\right).
\end{equation}

The cell size at level $l$ is $\Delta x_l = 2^{-l}$ (in units where
\texttt{boxlen} = 1).  The offset of child cell \texttt{ind} from the oct
centre is
\begin{equation}
  \boldsymbol{\delta}_\texttt{ind}
  = \left(i_x - \tfrac{1}{2},\; i_y - \tfrac{1}{2},\; i_z - \tfrac{1}{2}\right)
    \times \Delta x_l.
\end{equation}

The absolute position of child \texttt{ind} in grid $g$ in fractional box
units $[0, 1]$ is
\begin{equation}
  \tilde{x} = \frac{x_g + \delta_{\texttt{ind},x} - x_{\mathrm{bound},x}}{n_x}
             + \frac{1}{2},
\end{equation}
and similarly for $\tilde{y}$ and $\tilde{z}$ (the $+1/2$ shifts the origin
from the grid corner to the box corner).  In the code:
\begin{lstlisting}
xleaf = (xg(j,1) + xc(ind,1) - xbound(1)) / dble(nx) + 0.5d0
\end{lstlisting}

After reading, leaf positions are converted to physical units (cm):
\begin{equation}
  x_\mathrm{phys} = \tilde{x} \times L_\mathrm{box},
  \quad L_\mathrm{box} = \texttt{boxlen} \times \texttt{unit\_l}.
\end{equation}

%=============================================================================
\section{The Hydro File (\texttt{hydro\_NNNNN.outCCCCC})}
%=============================================================================

The hydro file stores the conserved hydrodynamic variables for every cell
(leaf and internal alike).  It is organised in strict lockstep with the AMR
file: one block per level per domain.

\subsection{Header Records}

\begin{center}
\begin{tabular}{rll}
\toprule
Record & Variable & Description \\
\midrule
1 & \texttt{ncpu}     & Number of MPI ranks \\
2 & \texttt{nvar}     & Number of hydro variables per cell \\
3 & \texttt{ndim}     & Number of dimensions \\
4 & \texttt{nlevelmax}& Maximum AMR level \\
5 & \texttt{nboundary}& Number of boundary domains \\
6 & \texttt{gamma}    & Adiabatic index \\
\bottomrule
\end{tabular}
\end{center}

\subsection{Variable Ordering}

Standard RAMSES hydro stores the following variables in order:

\begin{center}
\begin{tabular}{cll}
\toprule
Index & Variable & Physical meaning \\
\midrule
1 & $\rho$           & Gas mass density [code units] \\
2 & $\rho v_x$       & $x$-momentum density [code units] \\
3 & $\rho v_y$       & $y$-momentum density [code units] \\
4 & $\rho v_z$       & $z$-momentum density [code units] \\
5 & $E$              & Total energy density (kinetic + thermal) [code units] \\
6 & $Z/Z_\odot$      & Metallicity (if \texttt{metal=.true.}) \\
7$+$ & passive scalars & Additional tracers (simulation-dependent) \\
\bottomrule
\end{tabular}
\end{center}

\paragraph{Level block structure.}
For each level $l$ and domain $j$, the hydro file contains:
\begin{enumerate}
  \item Two header records (level number and domain number) — skipped.
  \item For each cell index \texttt{ind}$=1,\ldots,8$ and each variable
        \texttt{ivar}$=1,\ldots,\texttt{nvar}$: one record of length
        \texttt{ngridfile}$(j,l)$ containing the values for all grids of
        domain $j$ at level $l$.
\end{enumerate}
All records have the same element count \texttt{ngridfile}$(j,l)$, so the
loop structure is:

\begin{lstlisting}
do ind = 1, twotondim     ! 8 cells per oct
  do ivar = 1, nvar       ! variables
    read(iu_hyd) var(:, ind, ivar)  ! (ngrida) values
  end do
end do
\end{lstlisting}

\paragraph{Precision.}
RAMSES writes hydro variables in double precision (\texttt{real*8}) by
default.  Some configurations compile with single precision (\texttt{real*4}).
The LaRT reader handles both via the flag \texttt{ramses\_hydro\_prec}
(set to 8 or 4 in \texttt{read\_ramses\_amr.f90}).

%=============================================================================
\section{Physical Unit Conversions}
%=============================================================================

\subsection{Hydrogen Number Density}

The mass density in code units is converted to CGS and then to hydrogen
number density assuming a fully ionised gas with mean molecular weight
$\mu_\mathrm{H} \approx 1$ (one proton mass per hydrogen nucleus):
\begin{equation}
  n_H \;[\mathrm{cm}^{-3}]
  = \frac{\rho_\mathrm{code} \times \texttt{unit\_d}}{m_H},
  \quad m_H = 1.6726 \times 10^{-24}\;\mathrm{g}.
\end{equation}
The neutral fraction $n_\mathrm{HI}/n_H$ is computed separately in
\texttt{grid\_mod\_amr.f90} using collisional ionisation equilibrium (CIE) if
\texttt{par\%use\_cie\_condition = .true.}, otherwise all hydrogen is assumed
neutral.

\subsection{Velocity}

RAMSES stores momentum density $\rho v_\alpha$ (variable indices 2--4).
The velocity in CGS is:
\begin{equation}
  v_\alpha \;[\mathrm{cm\,s}^{-1}]
  = \frac{(\rho v_\alpha)_\mathrm{code}}{\rho_\mathrm{code}}
    \times \texttt{unit\_v},
  \quad \texttt{unit\_v} = \frac{\texttt{unit\_l}}{\texttt{unit\_t}}.
\end{equation}
LaRT receives velocities in km~s$^{-1}$ (divided by $10^5$).

\subsection{Temperature}

RAMSES variable 5 is the \emph{total energy density} $E$ (kinetic plus
thermal).  The thermal energy density is:
\begin{equation}
  e_\mathrm{th} = E - \frac{1}{2}\rho\,v^2
  = E - \frac{(\rho v_x)^2 + (\rho v_y)^2 + (\rho v_z)^2}{2\rho}.
\end{equation}
The specific thermal energy (per unit mass) in physical units is:
\begin{equation}
  \varepsilon \;[\mathrm{erg\,g}^{-1}]
  = \frac{e_\mathrm{th}}{\rho_\mathrm{code}}
    \times \texttt{unit\_v}^2.
\end{equation}
The temperature follows from the ideal gas law with adiabatic index $\gamma$
and mean molecular weight $\bar\mu$ (LaRT assumes $\bar\mu = 1.22$ for
neutral gas):
\begin{equation}
  T \;[\mathrm{K}]
  = (\gamma - 1)\;\varepsilon\;
    \frac{\bar\mu\,m_H}{k_B},
  \quad k_B = 1.381 \times 10^{-16}\;\mathrm{erg\,K}^{-1}.
\end{equation}
A floor of $T = 10\;\mathrm{K}$ is applied after conversion.

%=============================================================================
\section{The LaRT Reader: Step-by-Step}
%=============================================================================

The reader in \texttt{read\_ramses\_amr.f90} proceeds as follows:

\begin{enumerate}
  \item \textbf{Read \texttt{info} file.}  Extract \texttt{ncpu},
        \texttt{unit\_l}, \texttt{unit\_d}, \texttt{unit\_t},
        \texttt{boxlen}, and $\gamma$.

  \item \textbf{Count leaf cells (first pass).}
        \texttt{ramses\_count\_leaves} opens every
        \texttt{amr\_*.outCCCCC} file, reads the \texttt{son} arrays for
        all levels, and counts cells where $\texttt{son} = 0$.  This gives
        the total $N_\mathrm{leaf}$ needed to pre-allocate output arrays.

  \item \textbf{Read leaf data (second pass).}
        \texttt{ramses\_read\_all\_cpus} opens \emph{both} the AMR and
        hydro files simultaneously for each CPU:
        \begin{itemize}
          \item Reads AMR grids level by level, extracting oct centres
                \texttt{xg} and the \texttt{son} array.
          \item Reads hydro variables for the same grids.
          \item For each child cell with $\texttt{son} = 0$ (leaf): computes
                its physical position, $n_H$, $T$, and $\mathbf{v}$, and
                appends them to the output arrays.
        \end{itemize}

  \item \textbf{Convert positions.}  After the second pass, leaf positions
        (stored as fractions of \texttt{boxlen}) are multiplied by
        $L_\mathrm{cm} = \texttt{boxlen} \times \texttt{unit\_l}$ to give
        physical positions in cm.  These are then passed to
        \texttt{amr\_build\_tree} as-is (with \texttt{distance2cm = 1}).

  \item \textbf{Build octree and normalise.}  Handled by
        \texttt{grid\_mod\_amr.f90} (see the AMR description document).
\end{enumerate}

\subsection{Domain Skipping}

The AMR and hydro files interleave data from all $N_\mathrm{cpu}$ domains
and boundary domains at each level.  The LaRT reader only \emph{stores} data
for the domain that matches \texttt{icpu} (the loop variable); data from
other domains are skipped with \texttt{skip\_records}.  This reproduces the
original RAMSES read pattern used in post-processing tools such as
\textsc{pymses} and \textsc{yt}.

%=============================================================================
\section{File Naming Convention}
%=============================================================================

\begin{center}
\begin{tabular}{ll}
\toprule
File pattern & Content \\
\midrule
\texttt{output\_NNNNN/info\_NNNNN.txt}         & Global metadata \\
\texttt{output\_NNNNN/amr\_NNNNN.outCCCCC}     & AMR tree for CPU \texttt{CCCCC}\\
\texttt{output\_NNNNN/hydro\_NNNNN.outCCCCC}   & Hydro variables for CPU \texttt{CCCCC}\\
\texttt{output\_NNNNN/gravity\_NNNNN.outCCCCC} & Gravitational potential (optional)\\
\texttt{output\_NNNNN/part\_NNNNN.outCCCCC}    & Particle data (optional)\\
\bottomrule
\end{tabular}
\end{center}

where \texttt{NNNNN} is the zero-padded 5-digit snapshot index and
\texttt{CCCCC} is the zero-padded 5-digit CPU rank (1-based).

%=============================================================================
\section{LaRT Input for RAMSES Data}
%=============================================================================

\begin{lstlisting}[style=shellstyle]
&parameters
 par%use_amr_grid  = .true.
 par%amr_type      = 'ramses'
 par%amr_file      = '/path/to/simulation'  ! repository directory
 par%amr_snapnum   = 10                     ! snapshot number (NNNNN)
 par%no_photons    = 1e6
 par%taumax        = 1.0e7
 par%DGR           = 0.0
 par%spectral_type = 'voigt'
 par%source_geometry = 'point'
 par%xs_point      = <x_cm>   ! source position in cm
 par%ys_point      = <y_cm>
 par%zs_point      = <z_cm>
 par%nxfreq  = 121
 par%nxim    = 64
 par%nyim    = 64
/
\end{lstlisting}

\paragraph{Notes.}
\begin{itemize}
  \item \texttt{par\%amr\_file} is the \emph{repository directory} containing
        \texttt{output\_NNNNN/} (not the output directory itself).
  \item \texttt{par\%amr\_snapnum} selects the snapshot (the integer
        \texttt{NNNNN}).
  \item For RAMSES data, \texttt{par\%distance2cm} is automatically set to
        $1.0$ because the reader returns positions in cm
        (pre-multiplied by \texttt{unit\_l}).
        The \texttt{par\%distance\_unit} parameter is ignored.
  \item \texttt{par\%taumax} is specified as the line-centre optical depth
        from the box centre to the $+z$ boundary.
  \item Source position (\texttt{xs\_point} etc.)\ must be given in cm,
        matching the units of the RAMSES leaf positions after conversion.
  \item If \texttt{par\%use\_cie\_condition = .true.}, the neutral fraction
        $n_\mathrm{HI}/n_H$ is computed from the temperature via
        collisional ionisation equilibrium; otherwise all hydrogen is assumed
        neutral.
\end{itemize}

%=============================================================================
\section{Generic Text Format vs.\ RAMSES Binary}
%=============================================================================

For cases where RAMSES binary output is unavailable (e.g., data from other
simulation codes, or hand-constructed test grids), LaRT provides the generic
text format described separately.  The two formats differ in the following
ways:

\begin{center}
\begin{tabular}{lll}
\toprule
Property & RAMSES binary & Generic text \\
\midrule
File format       & Fortran unformatted binary & Plain text \\
Position unit     & cm (auto, from \texttt{unit\_l}) & Code units (set by \texttt{distance\_unit}) \\
\texttt{par\%amr\_type} & \texttt{'ramses'}          & \texttt{'generic'} \\
\texttt{par\%amr\_snapnum} & required              & not used \\
Multi-CPU support & yes (reads all \texttt{.outCCCCC} files) & no (single file) \\
Physical units    & full CGS conversion built in & user's responsibility \\
\bottomrule
\end{tabular}
\end{center}

The generic text file format is:
\begin{lstlisting}[style=shellstyle]
nleaf  boxlen
x  y  z  level  nH[cm^-3]  T[K]  vx[km/s]  vy[km/s]  vz[km/s]
...
\end{lstlisting}

The Python utility \texttt{examples/python/make\_amr\_grid.py} provides the
\texttt{AMRGrid} class to construct and write these files from arbitrary
density, temperature, and velocity fields.

%=============================================================================
\section{Known Limitations and Assumptions}
%=============================================================================

\begin{enumerate}
  \item \textbf{Neutral fraction.}  The RAMSES reader sets $n_H$ (total
        hydrogen density, not $n_\mathrm{HI}$).  The neutral fraction is
        computed later in \texttt{grid\_create\_amr} using either CIE
        (if \texttt{par\%use\_cie\_condition = .true.}) or the assumption
        of full neutrality.  For simulations with radiation feedback or
        on-the-fly ionisation, a custom neutral fraction field from a
        passive scalar may be needed.

  \item \textbf{Mean molecular weight.}  The temperature calculation in
        \texttt{ramses\_read\_all\_cpus} assumes $\bar\mu = 1.22$ (fully
        neutral primordial gas).  For ionised or metal-enriched gas this
        gives an overestimate.

  \item \textbf{Hydro precision.}  The variable \texttt{ramses\_hydro\_prec}
        (default~8) must be changed to~4 for simulations compiled with
        single-precision hydro.

  \item \textbf{Non-standard variable ordering.}  If the RAMSES simulation
        includes passive scalars before the energy variable (e.g., some
        MHD setups), the variable index for energy ($E =$ var~5) may be
        wrong.  Check the RAMSES namelist (\texttt{hydro.nml}) for the
        actual ordering.

  \item \textbf{Memory.}  All leaf cells are broadcast to every MPI rank
        in LaRT.  For large RAMSES outputs ($\gtrsim 10^7$ leaf cells)
        this may require significant RAM per node.
\end{enumerate}

%=============================================================================
\begin{thebibliography}{9}
\bibitem{teyssier2002}
Teyssier, R.\ 2002, A\&A, 385, 337
(\textit{RAMSES: a new high-resolution n-body and hydrodynamics code for
cosmological simulations})
\end{thebibliography}

\end{document}
